\documentclass[12pt,a4paper]{report}

\usepackage[utf8]{inputenc}
\usepackage[T1]{fontenc}
\usepackage[brazil]{babel}
\usepackage{geometry}
\usepackage{setspace}
\usepackage{amsmath,amssymb}
\usepackage{booktabs}
\usepackage{longtable}
\usepackage{array}
\usepackage{hyperref}
\usepackage{xcolor}
\usepackage{listings}

\geometry{left=3cm, right=2cm, top=3cm, bottom=2cm}
\onehalfspacing

\hypersetup{
    colorlinks=true,
    linkcolor=blue!60!black,
    urlcolor=blue!60!black,
    citecolor=blue!60!black,
    pdftitle={MiniLang-Compiler (AV2) --- Reconhecimento com Pilha e Fita Limitada},
    pdfauthor={Andrey Lourival Andrade Garcia; Igor Cecim Vilhena}
}

\begin{document}

\begin{titlepage}
    \centering
    {\large CENTRO UNIVERSITÁRIO DO ESTADO DO PARÁ --- CESUPA \par}
    {\normalsize Curso de Ciência da Computação --- Turma CC6NA \par}
    \vspace{2.5cm}
    {\LARGE \textbf{MiniLang-Compiler (AV2)} \par}
    \vspace{0.4cm}
    {\large \textbf{Reconhecimento de Linguagens em Duas Etapas com Autômato de Pilha (GLC) e Autômato Linearmente Limitado (LSC)} \par}
    \vspace{2.5cm}
    {\normalsize \textbf{Integrantes:} \par}
    {\normalsize Andrey Lourival Andrade Garcia \par}
    {\normalsize Igor Cecim Vilhena \par}
    \vspace{1.2cm}
    {\normalsize \textbf{Disciplina:} Linguagens Formais e Autômatos \par}
    {\normalsize \textbf{Docente:} Prof. Daniel Leal Souza \par}
    \vfill
    {\normalsize Belém --- PA \\ 25 de novembro de 2026 \par}
\end{titlepage}

\pagenumbering{roman}
\chapter*{Resumo}
\addcontentsline{toc}{chapter}{Resumo}

Este relatório técnico apresenta a modelagem formal, a implementação e a validação experimental do projeto \textbf{MiniLang-Compiler (AV2)}, evolução direta do analisador léxico \texttt{MiniLang-MiniLexer} desenvolvido na AV1. A aplicação valida cadeias formais de programas compactos da linguagem \texttt{MiniLang} em duas etapas sequenciais: (1) um \textbf{Autômato de Pilha Determinístico (AP)} com \textbf{9 estados} reconhece a estrutura sintática de escopos aninhados e seção de assinatura/chamada gerada por uma \textbf{Gramática Livre de Contexto (GLC)} LL(1) não-regular; e (2) caso o AP aceite, um \textbf{Autômato Linearmente Limitado Determinístico (ALL)} com \textbf{11 estados} opera sobre uma fita estritamente limitada de $n$ células graváveis e dois marcadores de borda fixos (\texttt{<} e \texttt{>}) para verificar a compatibilidade posicional de aridade e tipos entre a assinatura declarada e a chamada de função ($L_2 = \{w\texttt{:}w \mid w \in \{\texttt{i}, \texttt{s}\}^+\}$), uma dependência sensível ao contexto não reconhecível por GLCs. A solução inclui motor formal em Python, CLI, Interface Web Interativa com rastreamento animado da pilha e da fita, artefatos \texttt{.jff} para o JFLAP 7.1, 227 testes automatizados no \texttt{pytest} (135 preservados da AV1 + 92 da AV2) e validação direta no motor oficial do \texttt{JFLAP7.1.jar} com \textbf{0 divergências} em 418 casos.

\medskip
\noindent\textbf{Palavras-chave:} Linguagens Livres de Contexto; Gramática Livre de Contexto; Autômato de Pilha; Linguagens Sensíveis ao Contexto; Autômato Linearmente Limitado; MiniLang.

\tableofcontents
\clearpage
\pagenumbering{arabic}

\section{Problema, Motivação e Fundamentação Teórica}

\subsection{Contexto e Originalidade do Problema}

O projeto \textbf{MiniLang-Compiler (AV2)} é uma solução própria concebida como evolução direta do analisador léxico \texttt{MiniLang-MiniLexer} desenvolvido pela equipe na AV1~\cite{av1}. Enquanto a AV1 cobriu o nível \textbf{Tipo 3 (Regular)} da Hierarquia de Chomsky (Expressões Regulares e AFN$\varepsilon$ de Thompson), a AV2 implementa as duas etapas seguintes de um compilador real em um único pipeline de validação:

\begin{enumerate}
    \item \textbf{Etapa 1 --- Análise Sintática (\texttt{MiniParser}, Tipo 2 / LLC):} um Autômato de Pilha Determinístico (AP) valida a estrutura sintática de programas compactos da \texttt{MiniLang} contendo cabeçalho \texttt{p} (\texttt{program}), escopos \texttt{\{...\}} e expressões \texttt{(...)} aninhados e balanceados, seguidos do separador \texttt{\#}, de uma assinatura de parâmetros tipados (\texttt{i} = \texttt{int}, \texttt{s} = \texttt{string}), do separador \texttt{:}, dos argumentos da chamada e do terminador \texttt{;}.
    \item \textbf{Etapa 2 --- Análise Semântica (\texttt{MiniSemantic}, Tipo 1 / LSC):} executada \emph{exclusivamente} se o AP aceitar, um Autômato Linearmente Limitado Determinístico (ALL) recebe a fatia semântica \texttt{w:u} delimitada entre \texttt{\#} e \texttt{;} em uma fita finita com marcadores fixos \texttt{<} e \texttt{>}, verificando se a chamada possui exatamente a mesma aridade e a mesma ordem posicional de tipos declarada na assinatura ($w = u$).
\end{enumerate}

A escolha distancia-se integralmente de exercícios de sala e do exemplo da lauda~\cite{atividade} (parênteses e contagem de blocos \texttt{RRVVEE}), refletindo uma separação clássica documentada na literatura de compiladores~\cite{aho}: analisadores sintáticos baseados em GLC validam aninhamento arbitrário de blocos e parênteses, mas são teoricamente incapazes de impor igualdade posicional entre declaração e chamada de função.

\subsection{Classificação na Hierarquia de Chomsky}

\subsubsection{Por que $L_1$ é Livre de Contexto e não é Regular}

A linguagem sintática
\[
L_1 = \{ \texttt{p\{}\alpha\texttt{\}\#}w\texttt{:}u\texttt{;} \mid \alpha \in D_2^*, \, w, u \in \{\texttt{i}, \texttt{s}\}^+ \}
\]
é gerada pela GLC $G$ (Seção~2), logo $L_1 \in \text{LLC}$. Para provar que $L_1 \notin \text{LR}$, aplica-se o \textbf{Lema do Bombeamento para Linguagens Regulares}~\cite{sipser,hopcroft}:

\textit{Demonstração.} Suponha por contradição que $L_1$ seja regular com comprimento de bombeamento $p \ge 3$. Considere $z = \texttt{p\{(}^p\texttt{)}^p\texttt{\}\#i:i;} \in L_1$. Em qualquer decomposição $z = xyz$ com $|xy| \le p$ e $|y| \ge 1$, o fator $y$ encontra-se inteiramente contido no prefixo $\texttt{p\{(}^{p-2}$. Ao bombear $xy^2z$: (a) se $y$ contém \texttt{p}, surgem múltiplos cabeçalhos \texttt{p}; (b) se $y$ contém \texttt{\{}, cria-se excesso de chaves abertas sem fechamento; (c) se $y \in \texttt{(}^+$, o número de \texttt{(} excede o de \texttt{)}. Em todos os casos $xy^2z \notin L_1$. Contradição. Portanto, $L_1 \notin \text{LR}$. \hfill $\square$

\subsubsection{Por que $L_2$ é Sensível ao Contexto e não é LLC}

A linguagem semântica $L_2 = \{ w\texttt{:}w \mid w \in \{\texttt{i}, \texttt{s}\}^+ \}$ é reconhecida pelo ALL $M_{\text{ALL}}$ (Seção~3) em espaço estritamente linear ($n$ células graváveis + 2 bordas fixas), logo $L_2 \in \text{LSC}$~\cite{hopcroft}. Para provar que $L_2 \notin \text{LLC}$, aplica-se o \textbf{Lema do Bombeamento para LLCs}~\cite{sipser}:

\textit{Demonstração.} Suponha por contradição que $L_2$ seja livre de contexto com constante $p \ge 1$. Considere $z = \texttt{i}^p\texttt{s}^p\texttt{:}\texttt{i}^p\texttt{s}^p \in L_2$ e decomponha $z = uvwxy$ com $|vwx| \le p$ e $|vx| \ge 1$.
\begin{itemize}
    \item Se $vwx$ está inteiramente à esquerda (ou à direita) de \texttt{:}, então $uv^2wx^2y$ altera o comprimento de apenas uma metade, violando $w=u$.
    \item Se $vwx$ contém \texttt{:}, como $|vwx| \le p$, ele só alcança $\texttt{s}^p$ à esquerda e $\texttt{i}^p$ à direita. Se \texttt{:} ocorre em $v$ ou $x$, então $uv^0wx^0y$ remove \texttt{:} ($\notin L_2$). Caso contrário, $v \in \texttt{s}^*$ e $x \in \texttt{i}^*$, de modo que $uv^2wx^2y$ aumenta \texttt{s} à esquerda e/ou \texttt{i} à direita sem alterar \texttt{s} à direita nem \texttt{i} à esquerda, destruindo a igualdade posicional ($\notin L_2$).
\end{itemize}
Contradição. Logo, $L_2 \in \text{LSC} \setminus \text{LLC}$ e exige a memória bidirecional de um ALL. \hfill $\square$

\chapter{Etapa 1 --- Gramática Livre de Contexto e Autômato de Pilha}

\section{Especificação Formal da GLC ($G$)}

O alfabeto de terminais da Etapa 1 é $\Sigma_{\text{AP}} = \{\texttt{p}, \texttt{\{}, \texttt{\}}, \texttt{(}, \texttt{)}, \texttt{\#}, \texttt{i}, \texttt{s}, \texttt{:}, \texttt{;}\}$. A gramática $G = (V, T, P, S)$ é definida por $V = \{S, B, T, R\}$, $T = \Sigma_{\text{AP}}$, símbolo inicial $S$ e 9 produções LL(1) não-ambíguas:

\begin{align*}
P_1:\quad S &\to \texttt{p\{}\,B\,\texttt{\}\#}\,T\,\texttt{:}\,T\,\texttt{;} \\
P_2:\quad B &\to \texttt{\{}\,B\,\texttt{\}}\,B \qquad\qquad
P_3:\quad B \to \texttt{(}\,B\,\texttt{)}\,B \qquad\qquad
P_4:\quad B \to \varepsilon \\
P_5:\quad T &\to \texttt{i}\,R \qquad\qquad\qquad
P_6:\quad T \to \texttt{s}\,R \\
P_7:\quad R &\to \texttt{i}\,R \qquad\qquad\qquad
P_8:\quad R \to \texttt{s}\,R \qquad\qquad\qquad
P_9:\quad R \to \varepsilon
\end{align*}

\section{Derivação Passo a Passo de uma Cadeia Aceita ($\ge 18$ caracteres)}

Abaixo apresenta-se a derivação mais à esquerda completa da cadeia aceita \texttt{p\{(())\{\}\}\#isisi:isisi;} (22 caracteres, caso \texttt{AP-A01}), obtida de forma determinística pelo motor \texttt{src/av2/glc.py}:

\begin{enumerate}
    \item $S \Rightarrow \texttt{p\{}B\texttt{\}\#}T\texttt{:}T\texttt{;}$ \hfill $[P_1]$
    \item $\Rightarrow \texttt{p\{(}B\texttt{)}B\texttt{\}\#}T\texttt{:}T\texttt{;}$ \hfill $[P_3]$
    \item $\Rightarrow \texttt{p\{((}B\texttt{)}B\texttt{)}B\texttt{\}\#}T\texttt{:}T\texttt{;}$ \hfill $[P_3]$
    \item $\Rightarrow \texttt{p\{(())}B\texttt{)}B\texttt{\}\#}T\texttt{:}T\texttt{;}$ \hfill $[P_4]$
    \item $\Rightarrow \texttt{p\{(())}B\texttt{\}\#}T\texttt{:}T\texttt{;}$ \hfill $[P_4]$
    \item $\Rightarrow \texttt{p\{(())\{}B\texttt{\}}B\texttt{\}\#}T\texttt{:}T\texttt{;}$ \hfill $[P_2]$
    \item $\Rightarrow \texttt{p\{(())\{\}}B\texttt{\}\#}T\texttt{:}T\texttt{;}$ \hfill $[P_4]$
    \item $\Rightarrow \texttt{p\{(())\{\}\}\#}T\texttt{:}T\texttt{;}$ \hfill $[P_4]$
    \item $\Rightarrow \texttt{p\{(())\{\}\}\#i}R\texttt{:}T\texttt{;}$ \hfill $[P_5]$
    \item $\Rightarrow \texttt{p\{(())\{\}\}\#is}R\texttt{:}T\texttt{;}$ \hfill $[P_8]$
    \item $\Rightarrow \texttt{p\{(())\{\}\}\#isi}R\texttt{:}T\texttt{;}$ \hfill $[P_7]$
    \item $\Rightarrow \texttt{p\{(())\{\}\}\#isis}R\texttt{:}T\texttt{;}$ \hfill $[P_8]$
    \item $\Rightarrow \texttt{p\{(())\{\}\}\#isisi}R\texttt{:}T\texttt{;}$ \hfill $[P_7]$
    \item $\Rightarrow \texttt{p\{(())\{\}\}\#isisi:}T\texttt{;}$ \hfill $[P_9]$
    \item $\Rightarrow \texttt{p\{(())\{\}\}\#isisi:i}R\texttt{;}$ \hfill $[P_5]$
    \item $\Rightarrow \texttt{p\{(())\{\}\}\#isisi:is}R\texttt{;}$ \hfill $[P_8]$
    \item $\Rightarrow \texttt{p\{(())\{\}\}\#isisi:isi}R\texttt{;}$ \hfill $[P_7]$
    \item $\Rightarrow \texttt{p\{(())\{\}\}\#isisi:isis}R\texttt{;}$ \hfill $[P_8]$
    \item $\Rightarrow \texttt{p\{(())\{\}\}\#isisi:isisi}R\texttt{;}$ \hfill $[P_7]$
    \item $\Rightarrow \texttt{p\{(())\{\}\}\#isisi:isisi;}$ \hfill $[P_9]$
\end{enumerate}

\section{Especificação Formal do Autômato de Pilha ($M_{\text{AP}}$ --- 9 estados)}

Define-se $M_{\text{AP}} = (Q_{\text{AP}}, \Sigma_{\text{AP}}, \Gamma_{\text{AP}}, \delta_{\text{AP}}, q_0, Z, F_{\text{AP}})$ com:
\begin{itemize}
    \item \textbf{Estados (9 estados $\ge 8$):} $Q_{\text{AP}} = \{q_0, q_1, q_2, q_3, q_4, q_5, q_6, q_7, q_8\}$;
    \item \textbf{Alfabeto de entrada:} $\Sigma_{\text{AP}} = \{\texttt{p}, \texttt{\{}, \texttt{\}}, \texttt{(}, \texttt{)}, \texttt{\#}, \texttt{i}, \texttt{s}, \texttt{:}, \texttt{;}\}$;
    \item \textbf{Alfabeto da pilha:} $\Gamma_{\text{AP}} = \{Z, \texttt{\{}, \texttt{(}\}$ (com símbolo inicial $Z$);
    \item \textbf{Configuração inicial:} $(q_0, x, Z)$;
    \item \textbf{Critério de aceitação:} aceitação por \textbf{estado final} ($F_{\text{AP}} = \{q_8\}$) com \textbf{consumo integral da entrada} (e pilha final igual a $Z$).
\end{itemize}

\begin{table}[ht]
\centering
\caption{Função de transição $\delta_{\text{AP}}(q, a, X) = (q', \gamma)$ do DPDA de 9 estados.}
\small
\begin{tabular}{@{}lllll@{}}
\toprule
Estado & Lê & Desempilha & Destino & Empilha (topo à esq.) \\
\midrule
$q_0$ & \texttt{p} & $Z$ & $q_1$ & $Z$ \\
$q_1$ & \texttt{\{} & $Z$ & $q_2$ & \texttt{\{}$Z$ \\
$q_2$ & \texttt{\{} & \texttt{\{} / \texttt{(} & $q_2$ & \texttt{\{\{} / \texttt{\{(} \\
$q_2$ & \texttt{(} & \texttt{\{} / \texttt{(} & $q_2$ & \texttt{(\{} / \texttt{((} \\
$q_2$ & \texttt{\}} & \texttt{\{} & $q_2$ & $\varepsilon$ \\
$q_2$ & \texttt{)} & \texttt{(} & $q_2$ & $\varepsilon$ \\
$q_2$ & $\varepsilon$ & $Z$ & $q_3$ & $Z$ \\
$q_3$ & \texttt{\#} & $Z$ & $q_4$ & $Z$ \\
$q_4$ & \texttt{i} / \texttt{s} & $Z$ & $q_5$ & $Z$ \\
$q_5$ & \texttt{i} / \texttt{s} & $Z$ & $q_5$ & $Z$ \\
$q_5$ & \texttt{:} & $Z$ & $q_6$ & $Z$ \\
$q_6$ & \texttt{i} / \texttt{s} & $Z$ & $q_7$ & $Z$ \\
$q_7$ & \texttt{i} / \texttt{s} & $Z$ & $q_7$ & $Z$ \\
$q_7$ & \texttt{;} & $Z$ & $q_8$ & $Z$ \\
\bottomrule
\end{tabular}
\end{table}

\section{Justificativa da Equivalência $L(G) = L(M_{\text{AP}})$}

Os estados $q_0$ e $q_1$ impõem exatamente o prefixo \texttt{p\{} de $P_1$, empilhando \texttt{\{} sobre $Z$. Enquanto o escopo principal permanece aberto, a pilha em $q_2$ mantém o invariante de armazenar, acima de $Z$, exatamente a sequência de delimitadores \texttt{\{} e \texttt{(} abertos e ainda não fechados, o que corresponde biunivocamente às derivações de $B$ ($P_2, P_3, P_4$) até o fechamento do \texttt{\}} externo de $P_1$. Nesse exato instante o topo retorna a $Z$, habilitando exclusivamente a transição determinística $(q_2, \varepsilon, Z) \to (q_3, Z)$. Por fim, $q_3 \to q_4 \to q_5$ reconhece \texttt{\#} seguido de $\{\texttt{i},\texttt{s}\}^+$ (gerado por $T$ e $R$ via $P_5$--$P_9$), $q_5 \to q_6 \to q_7$ reconhece \texttt{:} seguido de $\{\texttt{i},\texttt{s}\}^+$ e $q_7 \to q_8$ consome \texttt{;}. Como $q_8$ não possui transições de saída, exige-se consumo integral da entrada. Logo $L(M_{\text{AP}}) = L(G)$.

\section{Etapa 2 --- LSC e Autômato Linearmente Limitado}

\subsection{Especificação Formal do ALL ($M_{\text{ALL}}$ --- 11 estados)}

Quando o AP aceita a entrada, o pipeline extrai a fatia semântica $x = w\texttt{:}u$ entre \texttt{\#} e \texttt{;} e a posiciona sobre a fita limitada do ALL:
\[
M_{\text{ALL}} = (Q_{\text{ALL}}, \Sigma_{\text{ALL}}, \Gamma_{\text{ALL}}, \texttt{<}, \texttt{>}, \delta_{\text{ALL}}, q_0, F_{\text{ALL}})
\]
com os seguintes componentes:
\begin{itemize}
    \item \textbf{Estados (11 $\ge 8$):} $Q_{\text{ALL}} = \{q_0, q_1, \dots, q_{10}\}$, final $F_{\text{ALL}} = \{q_{10}\}$;
    \item \textbf{Alfabetos:} $\Sigma_{\text{ALL}} = \{\texttt{i}, \texttt{s}, \texttt{:}\}$ e $\Gamma_{\text{ALL}} = \{\texttt{<}, \texttt{>}, \texttt{i}, \texttt{s}, \texttt{:}, \texttt{I}, \texttt{S}\}$;
    \item \textbf{Fita estritamente limitada:} para $|x| = n$, existem exatamente $n$ células graváveis (índices $1 \dots n$) e dois marcadores de borda fixos: \texttt{<} na célula $0$ e \texttt{>} na célula $n+1$;
    \item \textbf{Configuração inicial:} estado $q_0$, fita $\texttt{<}x\texttt{>}$, cabeçote na célula $0$ (compatível com o simulador de MT do JFLAP 7.1);
    \item \textbf{Decisão:} aceita em $q_{10}$ (sem transições de saída); rejeita ao parar em qualquer estado $q \notin F_{\text{ALL}}$ sem transição definida.
\end{itemize}

\begin{figure*}[t]
\centering
\resizebox{0.92\textwidth}{!}{%
\begin{tikzpicture}[
    >=Stealth,
    every state/.style={thick, minimum size=1.0cm, font=\small\ttfamily},
    every edge/.style={draw, ->, thick},
    lbl/.style={font=\scriptsize\ttfamily, align=center}
]
    \node[state, initial, initial text=] (q0) at (0, 0) {q0};
    \node[state] (q1) at (2.2, 0) {q1};
    \node[state] (q8) at (5.6, 0) {q8};
    \node[state] (q7) at (9.6, 0) {q7};
    \node[state] (q6) at (13.2, 0) {q6};

    \node[state] (q2) at (5.6, 2.5) {q2};
    \node[state] (q4) at (9.6, 2.5) {q4};

    \node[state] (q3) at (5.6, -2.5) {q3};
    \node[state] (q5) at (9.6, -2.5) {q5};

    \node[state] (q9) at (8.0, -4.9) {q9};
    \node[state, accepting] (q10) at (12.0, -4.9) {q10};

    \path
        (q0) edge node[above, lbl] {<$\to$<,R} (q1)
        (q1) edge node[above left, lbl] {i$\to$I,R} (q2)
        (q1) edge node[below left, lbl] {s$\to$S,R} (q3)

        (q2) edge[loop above, min distance=8mm] node[above, lbl] {i$\to$i,R \,|\, s$\to$s,R} (q2)
        (q2) edge node[above, lbl] {:$\to$:,R} (q4)
        (q4) edge[loop above, min distance=8mm] node[above, lbl] {I$\to$I,R \,|\, S$\to$S,R} (q4)
        (q4) edge node[above right, lbl] {i$\to$I,L} (q6)

        (q3) edge[loop below, min distance=8mm] node[below, lbl] {i$\to$i,R \,|\, s$\to$s,R} (q3)
        (q3) edge node[above, lbl] {:$\to$:,R} (q5)
        (q5) edge[loop below, min distance=8mm] node[below, lbl] {I$\to$I,R \,|\, S$\to$S,R} (q5)
        (q5) edge node[below right, lbl] {s$\to$S,L} (q6)

        (q6) edge[loop right, min distance=8mm] node[right, lbl] {I$\to$I,L \\ S$\to$S,L} (q6)
        (q6) edge node[above, lbl] {:$\to$:,L} (q7)
        (q7) edge[loop above, min distance=7mm] node[above, lbl] {i$\to$i,L \,|\, s$\to$s,L} (q7)
        (q7) edge node[above, lbl] {I$\to$I,R \,|\, S$\to$S,R} (q8)

        (q8) edge node[left, lbl] {i$\to$I,R} (q2)
        (q8) edge node[left, lbl] {s$\to$S,R} (q3)
        (q8) edge[bend left=18] node[right, pos=0.62, lbl] {:$\to$:,R} (q9)
        (q9) edge[loop below, min distance=8mm] node[below, lbl] {I$\to$I,R \,|\, S$\to$S,R} (q9)
        (q9) edge node[above, lbl] {>$\to$>,L} (q10);
\end{tikzpicture}%
}
\caption{Diagrama de estados do D-LBA $M_{\text{ALL}}$ (11 estados, equivalente a \texttt{docs/diagramas/AV2/all.jff}).}
\label{fig:all}
\end{figure*}

\begin{table*}[t]
\centering
\caption{Função de transição $\delta_{\text{ALL}}(q, X) = (q', Y, D)$ do D-LBA de 11 estados.}
\label{tab:all_delta}
\footnotesize
\setlength{\tabcolsep}{6pt}
\begin{tabular}{@{}llllll@{}}
\toprule
Estado & Lê & Destino & Escreve & Mov. & Papel semântico \\
\midrule
$q_0$ & \texttt{<} & $q_1$ & \texttt{<} & R & Valida borda esquerda e entra na célula 1 \\
$q_1$ & \texttt{i} / \texttt{s} & $q_2$ / $q_3$ & \texttt{I} / \texttt{S} & R & Exige assinatura não vazia e marca 1º tipo \\
$q_2$ & \texttt{i}, \texttt{s} / \texttt{:} & $q_2$ / $q_4$ & mesmo & R & Transporta \texttt{i} até o lado direito \\
$q_3$ & \texttt{i}, \texttt{s} / \texttt{:} & $q_3$ / $q_5$ & mesmo & R & Transporta \texttt{s} até o lado direito \\
$q_4$ & \texttt{I}, \texttt{S} / \texttt{i} & $q_4$ / $q_6$ & mesmo / \texttt{I} & R / L & Pula já casados e casa \texttt{i} $\to$ \texttt{I} \\
$q_5$ & \texttt{I}, \texttt{S} / \texttt{s} & $q_5$ / $q_6$ & mesmo / \texttt{S} & R / L & Pula já casados e casa \texttt{s} $\to$ \texttt{S} \\
$q_6$ & \texttt{I}, \texttt{S} / \texttt{:} & $q_6$ / $q_7$ & mesmo & L & Retorna à esquerda até o separador \texttt{:} \\
$q_7$ & \texttt{i}, \texttt{s} / \texttt{I}, \texttt{S} & $q_7$ / $q_8$ & mesmo & L / R & Retorna até a fronteira marcada e avança 1 \\
$q_8$ & \texttt{i} / \texttt{s} / \texttt{:} & $q_2$ / $q_3$ / $q_9$ & \texttt{I} / \texttt{S} / \texttt{:} & R & Marca próximo tipo ou inicia varredura final \\
$q_9$ & \texttt{I}, \texttt{S} / \texttt{>} & $q_9$ / $q_{10}$ & mesmo & R / L & Confirma ausência de sobras até \texttt{>} \\
\bottomrule
\end{tabular}
\end{table*}

\subsection{Garantia de Respeito aos Limites da Fita}

Em nenhuma transição do ALL o símbolo \texttt{<} ou \texttt{>} é sobrescrito, tampouco algum símbolo interno é substituído por marcador de borda. A única transição sobre \texttt{<} é $\delta(q_0, \texttt{<}) = (q_1, \texttt{<}, R)$ (move para dentro da fita) e a única transição sobre \texttt{>} é $\delta(q_9, \texttt{>}) = (q_{10}, \texttt{>}, L)$ (retorna para a célula $n$ e aceita). No retorno à esquerda em $q_7$, existe sempre ao menos um símbolo já marcado (\texttt{I} ou \texttt{S}) à direita de \texttt{<}, garantindo que o cabeçote reverta para $q_8$ antes de alcançar a borda esquerda. Tanto em \texttt{src/av2/all\_lba.py} quanto em \texttt{tests/JflapAv2RuntimeCheck.java}, asserções em tempo de execução confirmam a cada passo que $0 \le \text{cabeçote} \le n + 1$ e que \texttt{<} e \texttt{>} permanecem intactos.

\section{Implementação, Determinismo e Terminação}

\subsection{Arquitetura e Fonte Única de Verdade}

A aplicação foi implementada em \textbf{Python 3} (usando exclusivamente a biblioteca padrão para o motor formal, a CLI e a Interface Web Interativa, e \texttt{pytest} para a suíte automatizada de 92 testes da AV2), mantendo ativo o \texttt{MiniLexer} da AV1 em \texttt{src/} e arquivando os entregáveis legados da AV1 em \texttt{archive/av1/}.

O módulo \texttt{src/av2/spec.py} concentra, em estruturas imutáveis (\texttt{@dataclass(frozen=True, slots=True)}), as especificações canônicas da GLC, do AP e do ALL. A partir dele operam:
\begin{itemize}
    \item \texttt{src/av2/glc.py}: motor de derivação mais à esquerda LL(1);
    \item \texttt{src/av2/ap.py}: simulador formal de DPDA com pilha real (\texttt{list[str]}), validação de $\Sigma_{\text{AP}}$ e rastreamento completo;
    \item \texttt{src/av2/all\_lba.py}: simulador formal de D-LBA com fita fixa \texttt{['<', *entrada, '>']}, asserções de fronteira a cada passo e rastreamento completo;
    \item \texttt{src/av2/pipeline.py}: orquestração das duas etapas (o ALL só executa se o AP aceitar);
    \item \texttt{src/av2/minilang\_source.py}: ponte transparente via \texttt{MiniLexer} que projeta código \texttt{MiniLang} em $\Sigma_{\text{AP}}$ sem checagem sintática/semântica externa, mapeia \texttt{(linha, coluna)} para sublinhado de erro estilo IDE e sintetiza código para os 50 casos oficiais;
    \item \texttt{src/av2/cli.py} e \texttt{src/av2/web.py}: CLI e Web IDE estilo terminal (\texttt{./run}, \texttt{./ap}, \texttt{./all}, \texttt{./glc}, \texttt{./suite}, \texttt{./specs}, \texttt{./man}) com grafos SVG ao vivo, modo \textit{fullscreen}, pilha vertical, fita limitada e derivação;
    \item \texttt{src/av2/jflap\_io.py}: gerador e validador estrutural de \texttt{glc.jff}, \texttt{ap.jff} e \texttt{all.jff} compatíveis com o JFLAP 7.1~\cite{jflap}.
\end{itemize}

Em ambos os autômatos, a decisão resulta \textbf{estritamente} da consulta às tabelas $\delta_{\text{AP}}$ e $\delta_{\text{ALL}}$ e das operações reais de pilha e fita --- sem expressões regulares, contadores externos ou comparações diretas de strings.

\subsection{Determinismo e Ausência de Ramificação}

Tanto o AP quanto o ALL foram projetados como máquinas \textbf{100\% determinísticas}:
\begin{itemize}
    \item No \textbf{AP}, não existe par $(q, a, X)$ com mais de uma transição, e em $q_2$ nenhuma leitura atua sobre o topo $Z$; logo, a única $\varepsilon$-transição $(q_2, \varepsilon, Z) \to (q_3, Z)$ jamais conflita com leituras. O detector nativo \texttt{PDANondeterminismDetector} do JFLAP 7.1 confirmou \textbf{0 estados não-determinísticos} em \texttt{ap.jff}.
    \item No \textbf{ALL}, cada par $(q, X)$ mapeia para no máximo uma tripla $(q', Y, D)$, dispensando \textit{backtracking}: existe exatamente um caminho computacional por entrada e a rejeição ocorre imediatamente na ausência de transição.
\end{itemize}

\subsection{Prova de Terminação}

\begin{itemize}
    \item \textbf{Terminação do AP ($O(n)$):} toda transição de $\delta_{\text{AP}}$ avança a leitura em $+1$, exceto $(q_2, \varepsilon, Z) \to (q_3, Z)$. Como não existe retorno de $\{q_3, \dots, q_8\}$ para $q_2$, essa $\varepsilon$-transição ocorre no máximo uma vez por execução. Para entrada de tamanho $n$, o AP sempre para em no máximo $n + 1$ passos.
    \item \textbf{Terminação do ALL ($O(n^2)$):} cada iteração do ciclo principal ($q_1/q_8 \to q_2/q_3 \to q_4/q_5 \to q_6 \to q_7 \to q_8$) converte irreversivelmente um par de $\{\texttt{i}, \texttt{s}\}$ (um de cada lado de \texttt{:}) em $\{\texttt{I}, \texttt{S}\}$, sem nunca desmarcar células. Entre marcações, cada varredura move o cabeçote em direção monotônica sobre no máximo $n + 1$ células. Com no máximo $\lfloor n/2 \rfloor$ iterações, o ALL sempre para incondicionalmente em no máximo $(n+2)^2$ passos.
\end{itemize}

\section{Testes Obrigatórios e Análise dos Resultados}

Todos os 50 casos exigidos pela atividade~\cite{atividade} foram executados automaticamente via \texttt{pytest} (\textbf{92 testes automatizados da AV2 aprovados}; entregas históricas da AV1 arquivadas em \texttt{archive/av1/}), via CLI (\texttt{python -m src.av2.cli --suite}) e validados em runtime diretamente no simulador oficial do \texttt{JFLAP7.1.jar} (\textbf{418 casos, 0 divergências}). Além das Tabelas~\ref{tab:casos_ap}--\ref{tab:casos_int}, o ALL foi validado exaustivamente contra a especificação matemática $L_2$ em todas as \textbf{3.280 cadeias} de comprimento $0 \dots 7$ sobre $\{\texttt{i}, \texttt{s}, \texttt{:}\}$.

\begin{table*}[t]
\centering
\caption{Testes isolados do AP (10 aceitas + 10 rejeitadas + 2 casos-limite).}
\label{tab:casos_ap}
\scriptsize\setlength{\tabcolsep}{4pt}
\begin{tabularx}{\textwidth}{@{}llccX@{}}
\toprule
ID & Entrada & Esp. & Obt. & Regra verificada \\
\midrule
AP-A01 & \texttt{p\{(())\{\}\}\#isisi:isisi;} & Aceita & Aceita & Aninhamento misto de blocos e expressões com assinatura e chamada válidas (22 chars). \\
AP-A02 & \texttt{p\{((()))\}\#isisis:sisi;} & Aceita & Aceita & Parênteses profundamente aninhados; AP aceita mesmo com aridades diferentes (22 chars). \\
AP-A03 & \texttt{p\{\{(\{\})\}\}\#iiisss:ssii;} & Aceita & Aceita & Escopos aninhados em 3 níveis com expressão interna e tipos variados (22 chars). \\
AP-A04 & \texttt{p\{()()()\}\#sisisi:isii;} & Aceita & Aceita & Três expressões sequenciais no escopo principal (22 chars). \\
AP-A05 & \texttt{p\{\{\}()\{\}\}\#iissii:siss;} & Aceita & Aceita & Combinação sequencial de sub-blocos vazios e expressão vazia (22 chars). \\
AP-A06 & \texttt{p\{(\{()\})\}\#ssssss:iiii;} & Aceita & Aceita & Alternância estrita entre parênteses e chaves em 4 níveis de pilha (22 chars). \\
AP-A07 & \texttt{p\{(())\{\}\}\#isisiisis:isisiisis;} & Aceita & Aceita & Cadeia completa reaproveitada na integração (cenário aceito por ambos, 30 chars). \\
AP-A08 & \texttt{p\{(())\{\}\}\#isisiisis:isisiissi;} & Aceita & Aceita & Sintaxe correta com ordem de tipos trocada --- AP aceita (cenário 3 de integração, 30 chars). \\
AP-A09 & \texttt{p\{\{\{()\}\}\}\#is:isisiisis;} & Aceita & Aceita & Escopo triplo com assinatura curta e chamada longa (23 chars). \\
AP-A10 & \texttt{p\{(\{\})()\}\#iiiiiiii:ss;} & Aceita & Aceita & Escopo com bloco contendo expressão seguida de expressão irmã (22 chars). \\
AP-R01 & \texttt{p\{(())\{\}\#isisiisis:isisiisis;} & Rejeita & Rejeita & Chave '\{' do escopo principal não fechada antes de '\#' (cenário 2 de integração, 29 chars). \\
AP-R02 & \texttt{p\{(\{)\})\}\#isisis:isisis;} & Rejeita & Rejeita & Cruzamento inválido de delimitadores: ')' tentando fechar '\{' no topo da pilha (23 chars). \\
AP-R03 & \texttt{\{(())\{\}\}\#isisis:isisis;} & Rejeita & Rejeita & Ausência do cabeçalho obrigatório 'p' no início do programa (23 chars). \\
AP-R04 & \texttt{p\{(())\{\}\}\}\#isisi:isisi;} & Rejeita & Rejeita & Chave '\}' excedente após o fechamento do escopo principal (23 chars). \\
AP-R05 & \texttt{p\{(())\{\}\}isisi:isisis;} & Rejeita & Rejeita & Ausência do separador '\#' entre escopos e seção de assinatura/chamada (22 chars). \\
AP-R06 & \texttt{p\{(())\{\}\}\#:isisisisisi;} & Rejeita & Rejeita & Bloco de assinatura vazio antes de ':' (violação de T -> iR | sR, 23 chars). \\
AP-R07 & \texttt{p\{(())\{\}\}\#isisisisisi:;} & Rejeita & Rejeita & Bloco de chamada vazio após ':' (violação de T -> iR | sR, 23 chars). \\
AP-R08 & \texttt{p\{(())\{\}\}\#isisis:isisis} & Rejeita & Rejeita & Ausência do terminador ';' ao final da cadeia (estado final q8 não atingido, 23 chars). \\
AP-R09 & \texttt{p\{(())\{\}\}\#isisi:isisi;;} & Rejeita & Rejeita & Símbolo excedente ';' após o estado final q8 (entrada não pode ter sobras, 23 chars). \\
AP-R10 & \texttt{p\{(())\{\}\}\#is:is:isisis;} & Rejeita & Rejeita & Separador ':' duplicado na seção de assinatura/chamada (23 chars). \\
AP-L01 & \texttt{\textquotedbl\textquotedbl} & Rejeita & Rejeita & Caso-limite: cadeia vazia (rejeitada sem consumir entrada, permanece em q0). \\
AP-L02 & \texttt{p\{\}\#i:i;} & Aceita & Aceita & Caso-limite: menor cadeia aceita pela GLC e pelo AP (8 caracteres). \\
\bottomrule
\end{tabularx}
\end{table*}

\begin{table*}[t]
\centering
\caption{Testes isolados do ALL (10 aceitas + 10 rejeitadas + 3 casos-limite).}
\label{tab:casos_all}
\scriptsize\setlength{\tabcolsep}{4pt}
\begin{tabularx}{\textwidth}{@{}llccX@{}}
\toprule
ID & Entrada & Esp. & Obt. & Regra verificada \\
\midrule
ALL-A01 & \texttt{isisiisis:isisiisis} & Aceita & Aceita & Cópia exata de 9 parâmetros mistos (reaproveitada na integração, 19 chars). \\
ALL-A02 & \texttt{iiiiiiiii:iiiiiiiii} & Aceita & Aceita & Nove parâmetros homogêneos do tipo 'i' (int) idênticos (19 chars). \\
ALL-A03 & \texttt{sssssssss:sssssssss} & Aceita & Aceita & Nove parâmetros homogêneos do tipo 's' (string) idênticos (19 chars). \\
ALL-A04 & \texttt{iiiiisssss:iiiiisssss} & Aceita & Aceita & Dois blocos contíguos (5 'i' seguidos de 5 's') espelhados diretamente (21 chars). \\
ALL-A05 & \texttt{sssssiiiii:sssssiiiii} & Aceita & Aceita & Dois blocos contíguos (5 's' seguidos de 5 'i') idênticos (21 chars). \\
ALL-A06 & \texttt{isisisisis:isisisisis} & Aceita & Aceita & Alternância estrita de 10 parâmetros 'is' repetida integralmente (21 chars). \\
ALL-A07 & \texttt{siiisssiii:siiisssiii} & Aceita & Aceita & Padrão assimétrico de 10 tipos casado posição a posição (21 chars). \\
ALL-A08 & \texttt{iisssiiis:iisssiiis} & Aceita & Aceita & Assinatura assimétrica de 9 parâmetros casada exatamente (19 chars). \\
ALL-A09 & \texttt{sssiisssi:sssiisssi} & Aceita & Aceita & Sequência de 9 parâmetros começando por 's' e terminando em 'i' (19 chars). \\
ALL-A10 & \texttt{isssiiisss:isssiiisss} & Aceita & Aceita & Sequência simétrica de 10 parâmetros validada em fita limitada (21 chars). \\
ALL-R01 & \texttt{isisiisis:isisiissi} & Rejeita & Rejeita & Erro de ordem posicional nos 2 últimos tipos ('is' vs 'si'); mesmas quantidades (19 chars). \\
ALL-R02 & \texttt{isisiisis:isisiisisi} & Rejeita & Rejeita & Erro de aridade: chamada possui 1 argumento excedente após casar assinatura (20 chars). \\
ALL-R03 & \texttt{isisiisisi:isisiisis} & Rejeita & Rejeita & Erro de aridade: chamada possui 1 argumento a menos do que a assinatura (20 chars). \\
ALL-R04 & \texttt{iiiiiiiii:iiiiiiiis} & Rejeita & Rejeita & Erro de tipo na última posição ('i' esperado, 's' encontrado, 19 chars). \\
ALL-R05 & \texttt{sssssssss:issssssss} & Rejeita & Rejeita & Erro de tipo logo na 1ª posição ('s' esperado, 'i' encontrado, 19 chars). \\
ALL-R06 & \texttt{iiiiisssss:sssssiiiii} & Rejeita & Rejeita & Erro de ordem: exatamente 5 'i' e 5 's' de cada lado, porém invertidos (21 chars). \\
ALL-R07 & \texttt{isisiisisisisiisis} & Rejeita & Rejeita & Ausência do separador ':' na fita (cabeçote atinge '>' em q2/q3, 18 chars). \\
ALL-R08 & \texttt{isisi:isisi:isisii} & Rejeita & Rejeita & Presença de dois separadores ':' na fita (rejeitado na varredura final q9, 18 chars). \\
ALL-R09 & \texttt{:isisiisisisisiisi} & Rejeita & Rejeita & Assinatura vazia antes do separador ':' (rejeitado imediatamente em q1, 18 chars). \\
ALL-R10 & \texttt{isisiisisisisiisi:} & Rejeita & Rejeita & Chamada vazia após o separador ':' (rejeitado ao encontrar '>' em q4/q5, 18 chars). \\
ALL-L01 & \texttt{\textquotedbl\textquotedbl} & Rejeita & Rejeita & Caso-limite: cadeia vazia (fita '<>', rejeitada imediatamente em q1 sem cruzar bordas). \\
ALL-L02 & \texttt{:} & Rejeita & Rejeita & Caso-limite: apenas o separador ':' sem tipos em nenhum lado (fita '<:>'). \\
ALL-L03 & \texttt{i:i} & Aceita & Aceita & Caso-limite: menor cadeia aceita pelo ALL (3 células graváveis + 2 marcadores fixos). \\
\bottomrule
\end{tabularx}
\end{table*}

\begin{table*}[t]
\centering
\caption{Testes de integração do pipeline completo (AP $\to$ ALL).}
\label{tab:casos_int}
\scriptsize\setlength{\tabcolsep}{4pt}
\begin{tabularx}{\textwidth}{@{}llccX@{}}
\toprule
ID & Entrada & Esp. & Obt. (Parada) & Regra verificada \\
\midrule
INT-01 & \texttt{p\{(())\{\}\}\#isisiisis:isisiisis;} & Aceita & Aceita (CONCLUIDO) & Cenário 1: AP aceita a estrutura sintática e ALL aceita a assinatura/chamada idêntica. \\
INT-02 & \texttt{p\{(())\{\}\#isisiisis:isisiisis;} & Rejeita & Rejeita (AP) & Cenário 2: AP rejeita por '\{' não fechado; ALL não é executado. \\
INT-03 & \texttt{p\{(())\{\}\}\#isisiisis:isisiissi;} & Rejeita & Rejeita (ALL) & Cenário 3: AP aceita a sintaxe, mas ALL rejeita pela troca posicional dos tipos finais. \\
\bottomrule
\end{tabularx}
\end{table*}

\subsection{Testes com Símbolo Fora do Alfabeto (Erro Léxico)}
\begin{itemize}
    \item \textbf{LEX-01} --- Entrada: \texttt{p\{(())\{\}\}\#isisi@isisi;}. Saída: \texttt{Erro léxico (AP): símbolo inválido '@' na posição 15. Alfabeto aceito $\Sigma$\_AP = \{p, \{, \}, (, ), \#, i, s, :, ;\}.}
    \item \textbf{LEX-02} --- Entrada: \texttt{isisiisis\$isisiisis}. Saída: \texttt{Erro léxico (ALL): símbolo inválido '\$' na posição 9. Alfabeto aceito $\Sigma$\_ALL = \{i, s, :\}.}
\end{itemize}

\subsection{Validação Cruzada no JFLAP 7.1}
O script \texttt{tests/validate\_av2\_jflap\_runtime.py} executa \texttt{glc.jff}, \texttt{ap.jff} e \texttt{all.jff} diretamente sobre o \texttt{JFLAP7.1.jar}, confirmando:
\begin{itemize}
    \item \texttt{glc.jff}: 9 produções carregadas sem erro pelo \texttt{XMLCodec};
    \item \texttt{ap.jff}: 9 estados, \textbf{0 estados não-determinísticos} (via \texttt{PDANondeterminismDetector}), 34 casos, \textbf{0 divergências};
    \item \texttt{all.jff}: 11 estados, 384 casos (oficiais + enumeração completa até tamanho 5), com asserção de limites de fita no \texttt{TMSimulator}, \textbf{0 divergências}.
\end{itemize}

\subsection{Rastreamento Passo a Passo (Execuções Reais)}
As Tabelas~\ref{tab:trace_ap}, \ref{tab:trace_all_min} e \ref{tab:trace_all_int1} apresentam os rastreamentos reais produzidos pelos simuladores em Python sobre o cenário principal \texttt{INT-01} e o caso-limite \texttt{ALL-L03}, evidenciando as operações de pilha e os limites invioláveis \texttt{<} e \texttt{>} da fita.

\begin{table*}[t]
\centering
\caption{Rastreamento completo do AP sobre \texttt{INT-01} (\texttt{p\{(())\{\}\}\#isisiisis:isisiisis;}, 31 passos).}
\label{tab:trace_ap}
\scriptsize\setlength{\tabcolsep}{4pt}
\begin{tabularx}{\textwidth}{@{}rllllX@{}}
\toprule
Passo & Estado & Pos & Pilha & Restante & Transição \\
\midrule
0 & \texttt{q0} & 0 & \texttt{Z} & \texttt{p\{(())\{\}\}\#isisiisis:isisiisis;} & \texttt{(inicial)} \\
1 & \texttt{q1} & 1 & \texttt{Z} & \texttt{\{(())\{\}\}\#isisiisis:isisiisis;} & \texttt{$\delta$(q0, p, Z) = (q1, Z)} \\
2 & \texttt{q2} & 2 & \texttt{\{Z} & \texttt{(())\{\}\}\#isisiisis:isisiisis;} & \texttt{$\delta$(q1, \{, Z) = (q2, \{Z)} \\
3 & \texttt{q2} & 3 & \texttt{(\{Z} & \texttt{())\{\}\}\#isisiisis:isisiisis;} & \texttt{$\delta$(q2, (, \{) = (q2, (\{)} \\
4 & \texttt{q2} & 4 & \texttt{((\{Z} & \texttt{))\{\}\}\#isisiisis:isisiisis;} & \texttt{$\delta$(q2, (, () = (q2, (()} \\
5 & \texttt{q2} & 5 & \texttt{(\{Z} & \texttt{)\{\}\}\#isisiisis:isisiisis;} & \texttt{$\delta$(q2, ), () = (q2, $\varepsilon$)} \\
6 & \texttt{q2} & 6 & \texttt{\{Z} & \texttt{\{\}\}\#isisiisis:isisiisis;} & \texttt{$\delta$(q2, ), () = (q2, $\varepsilon$)} \\
7 & \texttt{q2} & 7 & \texttt{\{\{Z} & \texttt{\}\}\#isisiisis:isisiisis;} & \texttt{$\delta$(q2, \{, \{) = (q2, \{\{)} \\
8 & \texttt{q2} & 8 & \texttt{\{Z} & \texttt{\}\#isisiisis:isisiisis;} & \texttt{$\delta$(q2, \}, \{) = (q2, $\varepsilon$)} \\
9 & \texttt{q2} & 9 & \texttt{Z} & \texttt{\#isisiisis:isisiisis;} & \texttt{$\delta$(q2, \}, \{) = (q2, $\varepsilon$)} \\
10 & \texttt{q3} & 9 & \texttt{Z} & \texttt{\#isisiisis:isisiisis;} & \texttt{$\delta$(q2, $\varepsilon$, Z) = (q3, Z)} \\
11 & \texttt{q4} & 10 & \texttt{Z} & \texttt{isisiisis:isisiisis;} & \texttt{$\delta$(q3, \#, Z) = (q4, Z)} \\
12 & \texttt{q5} & 11 & \texttt{Z} & \texttt{sisiisis:isisiisis;} & \texttt{$\delta$(q4, i, Z) = (q5, Z)} \\
13 & \texttt{q5} & 12 & \texttt{Z} & \texttt{isiisis:isisiisis;} & \texttt{$\delta$(q5, s, Z) = (q5, Z)} \\
14 & \texttt{q5} & 13 & \texttt{Z} & \texttt{siisis:isisiisis;} & \texttt{$\delta$(q5, i, Z) = (q5, Z)} \\
15 & \texttt{q5} & 14 & \texttt{Z} & \texttt{iisis:isisiisis;} & \texttt{$\delta$(q5, s, Z) = (q5, Z)} \\
16 & \texttt{q5} & 15 & \texttt{Z} & \texttt{isis:isisiisis;} & \texttt{$\delta$(q5, i, Z) = (q5, Z)} \\
17 & \texttt{q5} & 16 & \texttt{Z} & \texttt{sis:isisiisis;} & \texttt{$\delta$(q5, i, Z) = (q5, Z)} \\
18 & \texttt{q5} & 17 & \texttt{Z} & \texttt{is:isisiisis;} & \texttt{$\delta$(q5, s, Z) = (q5, Z)} \\
19 & \texttt{q5} & 18 & \texttt{Z} & \texttt{s:isisiisis;} & \texttt{$\delta$(q5, i, Z) = (q5, Z)} \\
20 & \texttt{q5} & 19 & \texttt{Z} & \texttt{:isisiisis;} & \texttt{$\delta$(q5, s, Z) = (q5, Z)} \\
21 & \texttt{q6} & 20 & \texttt{Z} & \texttt{isisiisis;} & \texttt{$\delta$(q5, :, Z) = (q6, Z)} \\
22 & \texttt{q7} & 21 & \texttt{Z} & \texttt{sisiisis;} & \texttt{$\delta$(q6, i, Z) = (q7, Z)} \\
23 & \texttt{q7} & 22 & \texttt{Z} & \texttt{isiisis;} & \texttt{$\delta$(q7, s, Z) = (q7, Z)} \\
24 & \texttt{q7} & 23 & \texttt{Z} & \texttt{siisis;} & \texttt{$\delta$(q7, i, Z) = (q7, Z)} \\
25 & \texttt{q7} & 24 & \texttt{Z} & \texttt{iisis;} & \texttt{$\delta$(q7, s, Z) = (q7, Z)} \\
26 & \texttt{q7} & 25 & \texttt{Z} & \texttt{isis;} & \texttt{$\delta$(q7, i, Z) = (q7, Z)} \\
27 & \texttt{q7} & 26 & \texttt{Z} & \texttt{sis;} & \texttt{$\delta$(q7, i, Z) = (q7, Z)} \\
28 & \texttt{q7} & 27 & \texttt{Z} & \texttt{is;} & \texttt{$\delta$(q7, s, Z) = (q7, Z)} \\
29 & \texttt{q7} & 28 & \texttt{Z} & \texttt{s;} & \texttt{$\delta$(q7, i, Z) = (q7, Z)} \\
30 & \texttt{q7} & 29 & \texttt{Z} & \texttt{;} & \texttt{$\delta$(q7, s, Z) = (q7, Z)} \\
31 & \texttt{q8} & 30 & \texttt{Z} & \texttt{$\varepsilon$} & \texttt{$\delta$(q7, ;, Z) = (q8, Z)} \\
\bottomrule
\end{tabularx}
\end{table*}

\begin{table*}[t]
\centering
\caption{Rastreamento completo do ALL sobre o caso-limite \texttt{ALL-L03} (\texttt{i:i}, 9 passos).}
\label{tab:trace_all_min}
\scriptsize\setlength{\tabcolsep}{5pt}
\begin{tabularx}{\textwidth}{@{}rllllX@{}}
\toprule
Passo & Estado & Cabeçote & Lido & Fita (cabeçote em \texttt{[ ]}) & Transição \\
\midrule
0 & \texttt{q0} & 0 & \texttt{<} & \texttt{[<]i:i>} & \texttt{(inicial)} \\
1 & \texttt{q1} & 1 & \texttt{i} & \texttt{<[i]:i>} & \texttt{$\delta$(q0, <) = (q1, <, R)} \\
2 & \texttt{q2} & 2 & \texttt{:} & \texttt{<I[:]i>} & \texttt{$\delta$(q1, i) = (q2, I, R)} \\
3 & \texttt{q4} & 3 & \texttt{i} & \texttt{<I:[i]>} & \texttt{$\delta$(q2, :) = (q4, :, R)} \\
4 & \texttt{q6} & 2 & \texttt{:} & \texttt{<I[:]I>} & \texttt{$\delta$(q4, i) = (q6, I, L)} \\
5 & \texttt{q7} & 1 & \texttt{I} & \texttt{<[I]:I>} & \texttt{$\delta$(q6, :) = (q7, :, L)} \\
6 & \texttt{q8} & 2 & \texttt{:} & \texttt{<I[:]I>} & \texttt{$\delta$(q7, I) = (q8, I, R)} \\
7 & \texttt{q9} & 3 & \texttt{I} & \texttt{<I:[I]>} & \texttt{$\delta$(q8, :) = (q9, :, R)} \\
8 & \texttt{q9} & 4 & \texttt{>} & \texttt{<I:I[>]} & \texttt{$\delta$(q9, I) = (q9, I, R)} \\
9 & \texttt{q10} & 3 & \texttt{I} & \texttt{<I:[I]>} & \texttt{$\delta$(q9, >) = (q10, >, L)} \\
\bottomrule
\end{tabularx}
\end{table*}

\begin{table*}[t]
\centering
\caption{Amostra do rastreamento do ALL sobre \texttt{INT-01} (fatia \texttt{isisiisis:isisiisis}, 21 células, 201 passos: 1º ciclo e varredura final).}
\label{tab:trace_all_int1}
\scriptsize\setlength{\tabcolsep}{4pt}
\begin{tabularx}{\textwidth}{@{}rllllX@{}}
\toprule
Passo & Estado & Cab. & Lido & Fita (cabeçote em \texttt{[ ]}) & Transição \\
\midrule
0 & \texttt{q0} & 0 & \texttt{<} & \texttt{[<]isisiisis:isisiisis>} & \texttt{(inicial)} \\
1 & \texttt{q1} & 1 & \texttt{i} & \texttt{<[i]sisiisis:isisiisis>} & \texttt{$\delta$(q0, <) = (q1, <, R)} \\
2 & \texttt{q2} & 2 & \texttt{s} & \texttt{<I[s]isiisis:isisiisis>} & \texttt{$\delta$(q1, i) = (q2, I, R)} \\
3 & \texttt{q2} & 3 & \texttt{i} & \texttt{<Is[i]siisis:isisiisis>} & \texttt{$\delta$(q2, s) = (q2, s, R)} \\
4 & \texttt{q2} & 4 & \texttt{s} & \texttt{<Isi[s]iisis:isisiisis>} & \texttt{$\delta$(q2, i) = (q2, i, R)} \\
5 & \texttt{q2} & 5 & \texttt{i} & \texttt{<Isis[i]isis:isisiisis>} & \texttt{$\delta$(q2, s) = (q2, s, R)} \\
6 & \texttt{q2} & 6 & \texttt{i} & \texttt{<Isisi[i]sis:isisiisis>} & \texttt{$\delta$(q2, i) = (q2, i, R)} \\
7 & \texttt{q2} & 7 & \texttt{s} & \texttt{<Isisii[s]is:isisiisis>} & \texttt{$\delta$(q2, i) = (q2, i, R)} \\
8 & \texttt{q2} & 8 & \texttt{i} & \texttt{<Isisiis[i]s:isisiisis>} & \texttt{$\delta$(q2, s) = (q2, s, R)} \\
9 & \texttt{q2} & 9 & \texttt{s} & \texttt{<Isisiisi[s]:isisiisis>} & \texttt{$\delta$(q2, i) = (q2, i, R)} \\
10 & \texttt{q2} & 10 & \texttt{:} & \texttt{<Isisiisis[:]isisiisis>} & \texttt{$\delta$(q2, s) = (q2, s, R)} \\
11 & \texttt{q4} & 11 & \texttt{i} & \texttt{<Isisiisis:[i]sisiisis>} & \texttt{$\delta$(q2, :) = (q4, :, R)} \\
12 & \texttt{q6} & 10 & \texttt{:} & \texttt{<Isisiisis[:]Isisiisis>} & \texttt{$\delta$(q4, i) = (q6, I, L)} \\
13 & \texttt{q7} & 9 & \texttt{s} & \texttt{<Isisiisi[s]:Isisiisis>} & \texttt{$\delta$(q6, :) = (q7, :, L)} \\
14 & \texttt{q7} & 8 & \texttt{i} & \texttt{<Isisiis[i]s:Isisiisis>} & \texttt{$\delta$(q7, s) = (q7, s, L)} \\
15 & \texttt{q7} & 7 & \texttt{s} & \texttt{<Isisii[s]is:Isisiisis>} & \texttt{$\delta$(q7, i) = (q7, i, L)} \\
16 & \texttt{q7} & 6 & \texttt{i} & \texttt{<Isisi[i]sis:Isisiisis>} & \texttt{$\delta$(q7, s) = (q7, s, L)} \\
17 & \texttt{q7} & 5 & \texttt{i} & \texttt{<Isis[i]isis:Isisiisis>} & \texttt{$\delta$(q7, i) = (q7, i, L)} \\
18 & \texttt{q7} & 4 & \texttt{s} & \texttt{<Isi[s]iisis:Isisiisis>} & \texttt{$\delta$(q7, i) = (q7, i, L)} \\
19 & \texttt{q7} & 3 & \texttt{i} & \texttt{<Is[i]siisis:Isisiisis>} & \texttt{$\delta$(q7, s) = (q7, s, L)} \\
20 & \texttt{q7} & 2 & \texttt{s} & \texttt{<I[s]isiisis:Isisiisis>} & \texttt{$\delta$(q7, i) = (q7, i, L)} \\
21 & \texttt{q7} & 1 & \texttt{I} & \texttt{<[I]sisiisis:Isisiisis>} & \texttt{$\delta$(q7, s) = (q7, s, L)} \\
22 & \texttt{q8} & 2 & \texttt{s} & \texttt{<I[s]isiisis:Isisiisis>} & \texttt{$\delta$(q7, I) = (q8, I, R)} \\
\multicolumn{6}{c}{\textit{... ciclos intermediários de casamento (passos 23 a 188) ...}} \\
189 & \texttt{q7} & 9 & \texttt{S} & \texttt{<ISISIISI[S]:ISISIISIS>} & \texttt{$\delta$(q6, :) = (q7, :, L)} \\
190 & \texttt{q8} & 10 & \texttt{:} & \texttt{<ISISIISIS[:]ISISIISIS>} & \texttt{$\delta$(q7, S) = (q8, S, R)} \\
191 & \texttt{q9} & 11 & \texttt{I} & \texttt{<ISISIISIS:[I]SISIISIS>} & \texttt{$\delta$(q8, :) = (q9, :, R)} \\
192 & \texttt{q9} & 12 & \texttt{S} & \texttt{<ISISIISIS:I[S]ISIISIS>} & \texttt{$\delta$(q9, I) = (q9, I, R)} \\
193 & \texttt{q9} & 13 & \texttt{I} & \texttt{<ISISIISIS:IS[I]SIISIS>} & \texttt{$\delta$(q9, S) = (q9, S, R)} \\
194 & \texttt{q9} & 14 & \texttt{S} & \texttt{<ISISIISIS:ISI[S]IISIS>} & \texttt{$\delta$(q9, I) = (q9, I, R)} \\
195 & \texttt{q9} & 15 & \texttt{I} & \texttt{<ISISIISIS:ISIS[I]ISIS>} & \texttt{$\delta$(q9, S) = (q9, S, R)} \\
196 & \texttt{q9} & 16 & \texttt{I} & \texttt{<ISISIISIS:ISISI[I]SIS>} & \texttt{$\delta$(q9, I) = (q9, I, R)} \\
197 & \texttt{q9} & 17 & \texttt{S} & \texttt{<ISISIISIS:ISISII[S]IS>} & \texttt{$\delta$(q9, I) = (q9, I, R)} \\
198 & \texttt{q9} & 18 & \texttt{I} & \texttt{<ISISIISIS:ISISIIS[I]S>} & \texttt{$\delta$(q9, S) = (q9, S, R)} \\
199 & \texttt{q9} & 19 & \texttt{S} & \texttt{<ISISIISIS:ISISIISI[S]>} & \texttt{$\delta$(q9, I) = (q9, I, R)} \\
200 & \texttt{q9} & 20 & \texttt{>} & \texttt{<ISISIISIS:ISISIISIS[>]} & \texttt{$\delta$(q9, S) = (q9, S, R)} \\
201 & \texttt{q10} & 19 & \texttt{S} & \texttt{<ISISIISIS:ISISIISI[S]>} & \texttt{$\delta$(q9, >) = (q10, >, L)} \\
\bottomrule
\end{tabularx}
\end{table*}

\chapter{Limitações Observadas, Contribuições e Uso de IA}

\section{Limitações Observadas nos Testes}

Durante a execução dos testes e da validação em fita limitada, foram observadas três limitações práticas inerentes à modelagem adotada:

\begin{enumerate}
    \item \textbf{Complexidade quadrática do ALL em fita única ($O(n^2)$):} como o D-LBA percorre a fita de ida e volta a cada símbolo casado, cadeias longas crescem quadraticamente em número de passos (por exemplo, a cadeia de 19 caracteres \texttt{isisiisis:isisiisis} executa 201 passos, e uma assinatura de 50 parâmetros exige milhares de movimentos do cabeçote).
    \item \textbf{Subconjunto reduzido de tipos (\texttt{i} e \texttt{s}):} embora a \texttt{MiniLang} da AV1 possua cinco tipos (\texttt{int}, \texttt{float}, \texttt{string}, \texttt{bool}, \texttt{time}), a AV2 restringiu a seção semântica a dois tipos (\texttt{i} e \texttt{s}) para manter o diagrama do ALL limpo (11 estados) e verificável manualmente durante a arguição. Incluir os demais tipos exigiria adicionar um par de estados de transporte/casamento por novo tipo.
    \item \textbf{Assinatura única por cadeia:} o modelo valida exatamente uma declaração e uma chamada por entrada formal; suportar múltiplas funções declaradas com sobrecarga ou chamadas fora de ordem exigiria alfabeto auxiliar de identificadores e memória adicional na fita limitada.
\end{enumerate}

\section{Contribuição de Cada Integrante}

\begin{itemize}
    \item \textbf{Andrey Lourival Andrade Garcia:} evolução da especificação \texttt{MiniLang} (\texttt{MiniLang\_AV2\_Guia\_Projeto.md}); implementação do motor formal em Python (\texttt{spec.py}, \texttt{glc.py}, \texttt{ap.py}, \texttt{all\_lba.py}, \texttt{pipeline.py}); desenvolvimento da CLI e da Interface Web Interativa (\texttt{cli.py}, \texttt{web.py}, \texttt{static/}); geração e validação real dos arquivos \texttt{.jff} no \texttt{JFLAP7.1.jar}; suíte automatizada \texttt{pytest} e \texttt{README.md}.
    \item \textbf{Igor Cecim Vilhena:} pesquisa bibliográfica e fundamentação das provas na Hierarquia de Chomsky (Lema do Bombeamento para $L_1$ e $L_2$); revisão das derivações da GLC e dos 50 casos de teste da lauda; redação e revisão do relatório técnico (\texttt{docs/relatorio\_av2/}); elaboração dos slides (\texttt{docs/apresentacao\_av2/}) e análise das limitações observadas.
\end{itemize}

\section{Declaração de Uso de Inteligência Artificial}

Ferramentas de IA (\textbf{Google Antigravity / Gemini} na AV2; \textbf{Claude / ChatGPT} na base preservada da AV1) foram utilizadas como apoio na discussão de propostas originais de linguagem, revisão das demonstrações formais, estruturação do código em Python, da Interface Web Interativa, do validador Java do JFLAP e dos artefatos de documentação. A equipe verificou integralmente todas as transições de $\delta_{\text{AP}}$ e $\delta_{\text{ALL}}$, revisou as provas pelo Lema do Bombeamento e de terminação, executou os \textbf{227 testes no \texttt{pytest}} e validou os arquivos \texttt{glc.jff}, \texttt{ap.jff} e \texttt{all.jff} diretamente no \texttt{JFLAP7.1.jar} com \textbf{0 divergências}.

\section{Conclusão}

O projeto \texttt{MiniLang-Compiler (AV2)} cumpriu todos os requisitos técnicos, formais e experimentais da atividade, integrando análise léxica (AV1), sintática (GLC + DPDA de 9 estados) e semântica (LSC + D-LBA de 11 estados) em um fluxo coerente, determinístico, reproduzível e validado tanto em Python quanto no simulador oficial do JFLAP 7.1.


\renewcommand{\bibname}{Referências}
\begin{thebibliography}{9}
\addcontentsline{toc}{chapter}{Referências}

\bibitem{sipser}
SIPSER, Michael. \textit{Introduction to the Theory of Computation}. 3. ed. Boston: Cengage Learning, 2013.

\bibitem{hopcroft}
HOPCROFT, John E.; MOTWANI, Rajeev; ULLMAN, Jeffrey D. \textit{Introduction to Automata Theory, Languages, and Computation}. 3. ed. Boston: Pearson/Addison-Wesley, 2007.

\bibitem{aho}
AHO, Alfred V.; LAM, Monica S.; SETHI, Ravi; ULLMAN, Jeffrey D. \textit{Compilers: Principles, Techniques, and Tools}. 2. ed. Boston: Pearson/Addison-Wesley, 2006.

\bibitem{jflap}
RODGER, Susan H.; FINLEY, Thomas W. \textit{JFLAP: An Interactive Formal Languages and Automata Package}. Sudbury: Jones and Bartlett Publishers, 2006. Disponível em: \url{https://www.jflap.org/}.

\bibitem{atividade}
SOUZA, Daniel Leal. \textit{Atividade Prática AV2: Linguagens Livres de Contexto e Sensíveis ao Contexto}. Material da disciplina de Linguagens Formais e Autômatos, CESUPA, 2026.

\bibitem{av1}
GARCIA, Andrey Lourival Andrade; VILHENA, Igor Cecim. \textit{MiniLang-MiniLexer: Analisador Léxico da MiniLang (AV1)}. Repositório GitHub, 2026. Disponível em: \url{https://github.com/1drey2drey3drey/MiniLang-MiniLexer}.

\end{thebibliography}

\end{document}
